% Section 1.2, from lectures/L01/appendix/b_from_c_to_cpp.md.

\section{From C to Modern Embedded C++}
\label{c1:sec:cpp}\label{c1:app:b}

\subsection{Why use C++ in embedded systems?}
\label{c1:sec:why}
Embedded systems increasingly use C++ because it enables:
\begin{itemize}
  \item Better abstraction (via classes and interfaces).
  \item Stronger type safety through stricter compile-time checks.
  \item Compile-time guarantees (e.g.\ \code{static_assert}, templates).
  \item Improved modularity and separation of concerns.
  \item Reusable and testable drivers.
\end{itemize}
Importantly, C++ provides these benefits while still allowing low-level hardware access and
zero-cost abstractions.

Unlike many desktop C++ applications, embedded systems often:
\begin{itemize}
  \item Disable exceptions (and mark functions \code{noexcept}).
  \item Avoid dynamic allocation (\code{malloc/free} or \code{new/delete}), or limit it to system
    initialization.
  \item Prefer deterministic execution, meaning that the execution time of critical code paths must
    be predictable and bounded.
\end{itemize}

\subsection{Key differences between C and C++}
\label{c1:sec:differences}

\begin{narrowtable}{@{}lll@{}}
  \thead{Feature} & \thead{C} & \thead{C++} \\
  \midrule
  Namespaces & No & Yes \\
  Classes & No & Yes \\
  Templates & No & Yes \\
  References & No & Yes \\
  Compile-time polymorphism & No & Yes \\
  Standard containers & Limited & Rich STL \\
  RAII & No & Yes \\
\end{narrowtable}

\note RAII (\emph{Resource Acquisition Is Initialization}) means that a resource is acquired when
an object is created and automatically released when the object goes out of scope.

\subsection{\code{.h} vs \code{.hpp}}
\label{c1:sec:hpp}
In C and C++, header files are conventionally given the extension \code{.h}. In C++-only code,
the extension \code{.hpp} is also commonly used.

From the compiler's perspective, the extension makes no difference:
\begin{itemize}
  \item \code{#include} works identically regardless of the file extension.
  \item The compiler does not use the extension to decide whether to compile a file as C or C++;
    that is determined by the compiler invocation (e.g.\ \code{gcc} vs \code{g++}) or by the
    extension of the \emph{source} file that includes the header.
\end{itemize}
Even though the choice has no effect on the generated code, the extension is still useful as a
convention:
\begin{itemize}
  \item \code{.h} is shared between C and C++, so a \code{.h} file could be valid C, valid C++, or
    both --- the extension alone does not tell you which.
  \item \code{.hpp} immediately signals that a header uses C++-only features, such as classes,
    templates, namespaces, or references, and therefore requires a C++ compiler.
\end{itemize}
This is especially useful in a book like this one, where paired C and C++ implementations of the
same driver are sometimes shown side by side (see e.g.\ the \code{cpp_interface} example in
\secref{c3:sec:example}, and the C interface demo from the embedded C course that
Chapter~\ref{c3:ch:interfaces} links to). Using \code{.hpp} for the C++ version removes any doubt
about which language a given header belongs to.

\note In this book, header files use the \code{.hpp} extension.

\subsection{\code{extern "C"} and \code{\#ifdef __cplusplus}}
\label{c1:sec:externc}
C and C++ have different linkage rules. C++ supports function overloading, so the compiler
\emph{mangles} function names to encode parameter types, allowing multiple functions with the same
name but different signatures to coexist. C has no overloading and does not mangle names.

This becomes a problem when C++ code needs to call a function compiled by a C compiler (or vice
versa): if one side expects a mangled name and the other exports an unmangled one, the linker
cannot find a matching symbol.

\code{extern "C"} tells the C++ compiler to use C linkage for the enclosed declarations, disabling
name mangling for them:

\begin{cppcode}
extern "C" void gpioInit();
\end{cppcode}

\noindent Multiple declarations can be wrapped in a block:

\begin{cppcode}
extern "C"
{
void gpioInit();
void gpioWrite(std::uint8_t pin, bool state);
}
\end{cppcode}

\note[Important] \code{extern "C"} only affects linkage/name mangling. It does not make C++-only
features usable from C: a function with C linkage cannot be overloaded or be a template, and a
declaration that a C compiler is also going to read must be valid C (plain functions and data, no
classes, templates, overloading, or references).

\subsubsection{\code{\#ifdef __cplusplus}}
\code{extern "C"} is C++-only syntax; a C compiler does not understand it and will fail to compile
a header containing it. Since headers are often shared between C and C++ (see
\secref{c1:sec:hpp}), they are usually guarded with the predefined macro \code{__cplusplus},
which is only defined when the file is compiled by a C++ compiler:

\begin{ccode}
#ifdef __cplusplus
extern "C" {
#endif

void gpioInit();
void gpioWrite(uint8_t pin, bool state);

#ifdef __cplusplus
}
#endif
\end{ccode}

\begin{itemize}
  \item When compiled as C++, \code{__cplusplus} is defined, so the declarations are wrapped in
    \code{extern "C"}.
  \item When compiled as C, \code{__cplusplus} is undefined, the preprocessor strips out both
    \code{#ifdef} blocks entirely, and the C compiler only sees the plain declarations.
\end{itemize}
This pattern is common in embedded systems, where a hardware vendor typically ships C headers and
libraries (register definitions, HAL functions), while the application itself is written in C++.
Wrapping the vendor's C declarations this way lets C++ application code link against the C library
without modifying the vendor's code.

\note A header written this way is meant to be compiled as either C or C++, which is exactly the
kind of header that keeps the \code{.h} extension described in \secref{c1:sec:hpp} --- \code{.hpp}
would incorrectly imply it is C++-only.

\subsection{Namespaces}
\label{c1:sec:namespaces}
Namespaces are used to avoid symbol collisions in large software systems. They can be compared to
folders on a computer:
\begin{itemize}
  \item Just as different folders can contain files with the same name, different namespaces can
    contain functions, classes, variables, and other entities with the same name.
  \item For example, two different namespaces may both contain a function called \code{init()}
    without causing a conflict.
\end{itemize}
A namespace is defined using the \code{namespace} keyword, followed by a name and a block of code
enclosed in braces.

Below is a simple example of a namespace named \code{driver}:

\begin{cppcode}
namespace driver
{

} // namespace driver
\end{cppcode}

\noindent Anything declared inside this block becomes part of the \code{driver} namespace.

For example, the namespace may contain a function called \code{init()}:

\begin{cppcode}
namespace driver
{
void init()
{
    // Driver initialization code.
}
} // namespace driver
\end{cppcode}

\noindent To call this function from outside the namespace, we prefix it with the namespace name
followed by the scope resolution operator \code{::}:

\begin{cppcode}
int main()
{
    driver::init();
    return 0;
}
\end{cppcode}

\note
\begin{itemize}
  \item Most functionality from the C++ standard library is accessed through namespace
    \code{std}.
  \item Therefore, it is common to write \code{std::} before standard library types and functions,
    for example \code{std::uint8_t}, \code{std::uint16_t}, \code{std::size_t}, and
    \code{std::printf}.
  \item C library functionality is also available in C++, but it is usually included through C++
    headers such as \code{<cstdint>}, \code{<cstdio>}, and \code{<cstring>}.
  \item In modern C++, these headers are generally preferred over the older C-style forms
    \code{<stdint.h>}, \code{<stdio.h>}, and \code{<string.h>}.
\end{itemize}

\subsubsection{Anonymous namespaces}
In C++, an anonymous namespace can be used to restrict symbols (variables, functions, types, and
compile-time constants) to a single translation unit (i.e.\ a single source file).

This is similar to using the keyword \code{static} for functions or variables in C.

Anonymous namespaces are typically used in \code{.cpp} files to hide internal helper functions or
variables that should not be visible outside the file.

Example:

\begin{cppcode}
namespace
{
// Content only visible within this source file.
void helper() {}
} // namespace
\end{cppcode}

\subsubsection{Nested namespaces}
Namespaces can also be nested, similar to how folders can contain subfolders:
\begin{itemize}
  \item Namespaces are often used to reflect the logical structure of the project.
  \item For example, assume a project has the following directory structure:
\end{itemize}

\begin{console}
include/
    driver/
        gpio.hpp
        timer.hpp
source/
    driver/
        gpio.cpp
        timer.cpp
    main.cpp
\end{console}

\noindent A common convention would be:
\begin{itemize}
  \item The content of \file{driver/gpio.hpp} and \file{driver/gpio.cpp} belongs to the namespace
    \code{driver::gpio}.
  \item The content of \file{driver/timer.hpp} and \file{driver/timer.cpp} belongs to the namespace
    \code{driver::timer}.
\end{itemize}
This makes it easy to see which module a function belongs to.

\subsubsection{Nested namespace syntax}
Modern C++ allows nested namespaces to be declared in a compact form:

\begin{cppcode}
namespace driver::gpio
{
void init()
{
    // GPIO initialization code.
}
} // namespace driver::gpio
\end{cppcode}

\noindent Before C++17, the same namespace would typically be written as:

\begin{cppcode}
namespace driver
{
namespace gpio
{
void init()
{
    // GPIO initialization code.
}
} // namespace gpio
} // namespace driver
\end{cppcode}

\noindent Both forms are equivalent, but the compact form is preferred in modern C++.

\subsubsection{Using functions from nested namespaces}
A function defined in namespace \code{driver::gpio} can be called using its fully qualified name:

\begin{cppcode}
int main()
{
    driver::gpio::init();
    return 0;
}
\end{cppcode}

\noindent This makes it clear which module the function belongs to and prevents naming conflicts
in larger projects.

\note[In short] In embedded systems, namespaces are commonly used to organize drivers, hardware
abstraction layers (HAL), and communication protocols.

\subsection{\code{constexpr}}
\label{c1:sec:constexpr}
The keyword \code{constexpr} indicates that a value or expression (including simple functions) can
be evaluated at compile time rather than at runtime. This allows the compiler to compute values in
advance, which can improve performance and ensure that such values are usable in
constant-expression contexts.

In embedded systems, \code{constexpr} is particularly useful for defining hardware-related
constants such as clock frequencies, register offsets, buffer sizes, and bit masks.

For example:

\begin{cppcode}
/** CPU frequency in Hz. */
constexpr std::uint32_t cpuFrequency{16000000U};

/** Baud rate in bps (bits per second). */
constexpr std::uint32_t baudrate{115200U};

/** Buffer length in bytes. */
constexpr std::size_t bufLen{20U};
\end{cppcode}

\note Brace initialization \code{{}} is used instead of \code{=} above:
\begin{itemize}
  \item This is advantageous, since \code{{}} can be used to initialize everything.
  \item An exception is when using the \code{auto} keyword prior to C++17; then \code{auto} should
    be used in combination with \code{=} to avoid creating so-called initializer lists.
\end{itemize}

Because these values are known at compile time, the compiler can replace their usage directly with
the computed value in the generated machine code.

Compared to macros (\code{#define}), \code{constexpr} has several advantages:
\begin{itemize}
  \item It is type-safe.
  \item It obeys normal C++ scoping rules (functions, local blocks, and namespaces).
  \item It can be used in constant expressions.
\end{itemize}
For example, \code{constexpr} values can be used for:
\begin{itemize}
  \item Array sizes.
  \item Template parameters.
  \item Compile-time calculations.
\end{itemize}
In modern C++, \code{constexpr} is generally preferred over \code{#define} for defining constants.

\code{constexpr} can be used for functions when the values are known at compile time:

\begin{cppcode}
constexpr int add(const int x, const int y) { return x + y; }
\end{cppcode}

\subsection{\code{noexcept}}
\label{c1:sec:noexcept}
The keyword \code{noexcept} indicates that a function is guaranteed not to throw exceptions.

For reference, an example of throwing an exception is shown below. In this example, an exception is
thrown if the value of \code{val} exceeds \code{100U}:

\begin{cppcode}
#include <cstdint>
#include <stdexcept>

int main()
{
    constexpr std::uint8_t limit{100U};
    std::uint8_t val{};

    while (1)
    {
        if (limit < val++)
        {
            throw std::out_of_range("Value exceeded limit!");
        }
    }
    return 0;
}
\end{cppcode}

\noindent Exceptions are a mechanism used in many C++ programs to signal runtime errors:
\begin{itemize}
  \item When an exception is thrown, the program begins stack unwinding, meaning that the call stack
    is traversed until a matching exception handler (i.e.\ a \code{catch} block) is found.
  \item If a function marked \code{noexcept} attempts to throw an exception, the program will
    terminate immediately.
\end{itemize}
In many embedded systems, exceptions are disabled entirely. This is typically done to:
\begin{itemize}
  \item Reduce binary size.
  \item Simplify the runtime environment.
  \item Ensure deterministic behavior.
\end{itemize}
In such systems, errors are typically handled using C-style mechanisms such as:
\begin{itemize}
  \item Return codes.
  \item Status flags.
  \item Error callbacks.
\end{itemize}
Because of this, it is common practice in embedded C++ to mark functions that are not expected to
fail with the \code{noexcept} keyword.

Doing so has several benefits:
\begin{itemize}
  \item It documents that the function cannot throw exceptions.
  \item It enables additional compiler optimizations.
  \item It avoids the need for exception-related stack handling.
\end{itemize}
Example:

\begin{cppcode}
namespace driver::gpio
{
bool read() noexcept { return false; }
} // namespace driver::gpio
\end{cppcode}

\note[Important]
\begin{itemize}
  \item In embedded drivers, functions that perform simple hardware operations (such as reading or
    writing registers) are typically marked \code{noexcept}, since they are not expected to throw
    exceptions.
  \item Functions that may perform dynamic memory allocation, such as resizing a vector, should not
    be marked \code{noexcept}, since memory allocation may fail.
\end{itemize}

\subsection{Default arguments}
\label{c1:sec:defaultargs}
C++ allows functions to specify default values for parameters:
\begin{itemize}
  \item If the caller does not provide a value for a parameter, the default value is used
    automatically.
  \item This can simplify APIs by allowing common cases to use fewer arguments.
\end{itemize}
For example, consider the function \code{log()} below:

\begin{cppcode}
namespace debug
{
void log(const char* message, std::uint8_t level = 0U)
{
    // Log message with the specified level.
}
} // namespace debug
\end{cppcode}

\noindent The parameter \code{level} has a default value of \code{0U}. This means the function can
be called in two different ways:
\begin{itemize}
  \item If the log level is not important, the function can be called with only the message:
\begin{cppcode}
debug::log("System started");
\end{cppcode}
    In this case, the value \code{0U} will automatically be used for \code{level}.
  \item If a specific log level is desired, the caller can provide it explicitly:
\begin{cppcode}
debug::log("Driver failure", 2U);
\end{cppcode}
\end{itemize}
Default arguments are commonly used in embedded software for functions that have sensible default
behavior, such as:
\begin{itemize}
  \item Logging systems.
  \item Driver configuration functions.
  \item Initialization routines.
\end{itemize}
They allow APIs to remain flexible while keeping common usage simple.

\subsubsection{Default argument placement}
In C++, default arguments may only be omitted from right to left. This means that parameters with
default values must be placed at the end of the parameter list.

A valid example is shown below:

\begin{cppcode}
void printValue(const std::uint32_t value, const bool hex = false) noexcept;
\end{cppcode}

\noindent This function can be called in either of the following ways:

\begin{cppcode}
printValue(42U);
printValue(42U, true);
\end{cppcode}

\noindent The following is not valid:

\begin{cppcode}
void printValue(const std::uint32_t value = 0U, const bool hex) noexcept;
\end{cppcode}

\noindent This is not allowed because the second parameter does not have a default value, even
though the first one does.

If this were allowed, a call such as this:

\begin{cppcode}
printValue(true);
\end{cppcode}

\noindent would become unclear and potentially cause errors (\code{true} can be implicitly
converted to \code{value} as \code{1}).

\subsection{Modern C++ structs}
\label{c1:sec:structs}
In C++, a \code{struct} can contain not only data members but also member functions, often
referred to as methods. This allows related data and operations to be grouped together in a
single type.

Consider the following \code{driver::Gpio} driver, implemented as a stub (for simulation):

\begin{cppcode}
#include <cstdint>

namespace driver
{
/**
 * @brief GPIO stub driver.
 */
struct Gpio
{
    /** Pin the GPIO is connected to. */
    const std::uint8_t pin;

    /** GPIO state. */
    bool state;

    /**
     * @brief Set GPIO state.
     *
     * @param[in] state GPIO state (true = enabled, false = disabled).
     */
    void write(const bool state) noexcept
    {
        // Note: The 'this' keyword is used here to refer to our member variable 'state',
        // since we have an input argument with the same name.
        this->state = state;
    }

    /**
     * @brief Get GPIO state.
     *
     * @return True if the GPIO is enabled, false otherwise.
     *
     * @note The 'const' keyword is used after the method name to set the GPIO instance to
     *       read-only in the scope of this method.
     */
    bool read() const noexcept { return state; }
};
} // namespace driver
\end{cppcode}

\noindent The \code{driver::Gpio} struct contains:
\begin{itemize}
  \item Two data members (\code{pin} and \code{state}).
  \item Two methods (\code{write} and \code{read}).
\end{itemize}
These functions operate on the data stored in the struct.

Instances of the struct can be created as follows:

\begin{cppcode}
// Initialize LED connected to pin 9, state set to false.
driver::Gpio led{9U, false};

// Initialize button connected to pin 13, simulate pressdown at startup.
driver::Gpio button{13U, true};
\end{cppcode}

\noindent Note that unlike in C, the keyword \code{struct} is not required when creating an
instance of a struct and can therefore be omitted.

The functions can then be called using the dot operator, just like the data members:

\begin{cppcode}
// Enable the LED if the button is pressed.
const bool buttonPressed{button.read()};
led.write(buttonPressed);
\end{cppcode}

\noindent In traditional C, similar functionality would typically be implemented using a struct
combined with separate functions that take a pointer to the struct, as shown below:

\begin{ccode}
#include <stdbool.h>
#include <stddef.h>
#include <stdint.h>

/**
 * @brief GPIO stub driver.
 */
typedef struct
{
    /** Pin the GPIO is connected to. */
    const uint8_t pin;

    /** GPIO state. */
    bool state;
} gpio_t;

/**
 * @brief Set GPIO state.
 *
 * @param[in, out] self Pointer to the GPIO.
 * @param[in] state GPIO state (true = enabled, false = disabled).
 */
void gpio_write(gpio_t* self, const bool state)
{
    if (NULL == self) { return; }
    self->state = state;
}

/**
 * @brief Get GPIO state.
 *
 * @param[in] self Pointer to the GPIO.
 *
 * @return True if the GPIO is enabled, false otherwise.
 */
bool gpio_read(const gpio_t* self)
{
    return NULL != self ? self->state : false;
}
\end{ccode}

\noindent Instances of the struct can then be created and used as shown below:

\begin{ccode}
// Initialize LED connected to pin 9, state set to false.
gpio_t led = {9U, false};

// Initialize button connected to pin 13, simulate pressdown at startup.
gpio_t button = {13U, true};

// Enable the LED if the button is pressed.
const bool button_pressed = gpio_read(&button);
gpio_write(&led, button_pressed);
\end{ccode}

\subsubsection{Comparison}
\paragraph{C-style approach}

\begin{ccode}
gpio_write(&led, true);
\end{ccode}

\paragraph{C++ approach}

\begin{cppcode}
led.write(true);
\end{cppcode}

\noindent The C++ syntax improves readability and keeps related functionality together, making the
code easier to understand and maintain.

\subsubsection{Constructor and destructor}
It is also possible to add initialization and cleanup methods. These methods are known as:
\begin{itemize}
  \item Constructor:
  \begin{itemize}
    \item Has the same name as the struct.
    \item Called automatically when an object is created.
    \item Can be used to automatically initialize an object on creation (for instance pin
      configuration).
    \item For example, a constructor for the \code{driver::Gpio} struct above can be implemented
      as follows:
  \end{itemize}
\end{itemize}

\begin{cppcode}
/**
 * @brief Constructor.
 *
 * @param[in] pin Pin the GPIO is connected to.
 * @param[in] initialState Initial state (default = disabled).
 */
Gpio(const std::uint8_t pin, const bool initialState = false) noexcept
    : pin{pin}
    , state{initialState}
{
    std::printf("Initializing GPIO at pin %u!\n", pin);
    // Configure pin here!
}
\end{cppcode}

\begin{itemize}
  \item Destructor:
  \begin{itemize}
    \item Has the same name as the struct, but with prefix \code{~}.
    \item Called automatically when an object is destroyed.
    \item For example, this happens when an object goes out of scope or when \code{delete} is used.
    \item Can be used to automatically release allocated resources or undo hardware-related setup.
    \item Often used to release peripherals, disable interrupts, or restore hardware state in
      embedded systems.
  \end{itemize}
\end{itemize}

\begin{cppcode}
/**
 * @brief Destructor.
 */
~Gpio() noexcept
{
    std::printf("Releasing resources reserved for GPIO at pin %u!\n", pin);
    // Release allocated resources here.
}
\end{cppcode}

\noindent After adding the constructor and destructor, the struct \code{driver::Gpio} looks like
this:

\begin{cppcode}
#include <cstdint>
#include <cstdio>

namespace driver
{
/**
 * @brief GPIO stub driver.
 */
struct Gpio
{
    /** Pin the GPIO is connected to. */
    const std::uint8_t pin;

    /** GPIO state. */
    bool state;

    /**
     * @brief Constructor.
     *
     * @param[in] pin Pin the GPIO is connected to.
     * @param[in] initialState Initial state (default = disabled).
     */
    Gpio(const std::uint8_t pin, const bool initialState = false) noexcept
        : pin{pin}
        , state{initialState}
    {
        std::printf("Initializing GPIO at pin %u!\n", pin);
        // Configure pin here!
    }

    /**
     * @brief Destructor.
     */
    ~Gpio() noexcept
    {
        std::printf("Releasing resources reserved for GPIO at pin %u!\n", pin);
        // Release allocated resources here.
    }

    /**
     * @brief Set GPIO state.
     *
     * @param[in] state GPIO state (true = enabled, false = disabled).
     */
    void write(const bool state) noexcept
    {
        // Note: The 'this' keyword refers to the current object.
        // It is used here because the parameter name 'state' hides the member variable.
        this->state = state;
    }

    /**
     * @brief Get GPIO state.
     *
     * @return True if the GPIO is enabled, false otherwise.
     *
     * @note The 'const' keyword is used after the method name to set the GPIO instance to
     *       read-only in the scope of this method.
     */
    bool read() const noexcept { return state; }
};
} // namespace driver
\end{cppcode}

\noindent Before a constructor is defined, the struct can be initialized using aggregate
initialization. After a user-defined constructor is added, the same brace syntax may still be used,
but it now calls the constructor instead.

For example, the following initialization:

\begin{cppcode}
int main()
{
    driver::Gpio led{9U, false};
    led.write(true);
}
\end{cppcode}

\noindent will generate the following output because the constructor and destructor are called
automatically:

\begin{console}
Initializing GPIO at pin 9!
Releasing resources reserved for GPIO at pin 9!
\end{console}

\noindent This illustrates the idea behind \textbf{RAII} (\emph{Resource Acquisition Is
Initialization}):
\begin{itemize}
  \item Initialization is performed during object creation.
  \item Cleanup is performed automatically when the object is destroyed.
\end{itemize}

\subsubsection{Introduction to encapsulation}
In C++, parts of a struct can be made private using the \code{private} keyword, as shown below for
the \code{driver::Gpio} struct:
\begin{itemize}
  \item The \code{public} keyword has been added for clarity; in a C++ struct, everything is public
    by default.
  \item The \code{private} keyword has been added to create a private segment.
  \item The member variables and methods have been removed for simplicity.
\end{itemize}

\begin{cppcode}
namespace driver
{
struct Gpio
{
public:
    // Public segment - accessible outside the struct.
private:
    // Private segment - inaccessible outside the struct.
};
} // namespace driver
\end{cppcode}

\noindent Using the \code{private} keyword has an important advantage:
\begin{itemize}
  \item Symbols (member variables, constants, and methods) declared in the private segment are
    inaccessible outside the struct.
  \item This allows us to hide implementation details and data that should not be accessed directly
    by other parts of the program.
  \item In many designs, member variables are declared private. This ensures that the internal
    state of the object can only be modified through its member functions.
\end{itemize}
For instance, the member variables of the \code{driver::Gpio} struct are typically made private,
which is shown below:
\begin{itemize}
  \item The private segment is typically placed at the bottom of the struct, while the public
    interface (constructors and methods) is placed first. This makes it easier to see how the
    struct should be used.
  \item Here we add the \code{my} prefix to make it clear that these are member variables and to
    reduce the need to use the \code{this} keyword (\code{this->state = state} was used earlier).
\end{itemize}

\begin{cppcode}
namespace driver
{
struct Gpio
{
private:
    /** Pin the GPIO is connected to. */
    const std::uint8_t myPin;

    /** GPIO state. */
    bool myState;
};
} // namespace driver
\end{cppcode}

\note The snippet above only shows the private segment of \code{driver::Gpio}.

Then we update the struct:
\begin{itemize}
  \item Member variables \code{pin} and \code{state} are replaced with \code{myPin} and
    \code{myState}.
  \item The \code{pin()} method is added so that users still can read the pin number.
\end{itemize}

\begin{cppcode}
#include <cstdint>
#include <cstdio>

namespace driver
{
/**
 * @brief GPIO stub driver.
 */
struct Gpio
{
    /**
     * @brief Constructor.
     *
     * @param[in] pin Pin the GPIO is connected to.
     * @param[in] initialState Initial state (default = disabled).
     */
    Gpio(const std::uint8_t pin, const bool initialState = false) noexcept
        : myPin{pin}
        , myState{initialState}
    {
        std::printf("Initializing GPIO at pin %u!\n", myPin);
        // Configure pin here!
    }

    /**
     * @brief Destructor.
     */
    ~Gpio() noexcept
    {
        std::printf("Releasing resources reserved for GPIO at pin %u!\n", myPin);
        // Release allocated resources here.
    }

    /**
     * @brief Get GPIO pin number.
     *
     * @return Pin the GPIO is connected to.
     */
    std::uint8_t pin() const noexcept { return myPin; }

    /**
     * @brief Set GPIO state.
     *
     * @param[in] state GPIO state (true = enabled, false = disabled).
     */
    void write(const bool state) noexcept { myState = state; }

    /**
     * @brief Get GPIO state.
     *
     * @return True if the GPIO is enabled, false otherwise.
     */
    bool read() const noexcept { return myState; }

private:
    /** Pin the GPIO is connected to. */
    const std::uint8_t myPin;

    /** GPIO state. */
    bool myState;
};
} // namespace driver
\end{cppcode}

\noindent After making the member variables private, the compiler will generate an error if we try
to access them directly:

\begin{cppcode}
int main()
{
    driver::Gpio led{9U};

    // Allowed: pin() provides controlled read access to the pin number.
    std::printf("The LED is connected to pin %u!\n", led.pin());

    // Allowed: write() is part of the public interface.
    led.write(true);

    // Not allowed: myPin is private.
    // std::printf("Pin: %u!\n", led.myPin);

    // Not allowed: myState is private.
    // led.myState = true;
}
\end{cppcode}

\noindent Chapter~\ref{c2:ch:classes} introduces classes:
\begin{itemize}
  \item In C++, structs and classes support almost the same language features, but classes use
    \code{private} as the default access level, while structs use \code{public}.
  \item Structs are often used for simple data-oriented types, while classes are commonly used for
    more complex objects that may involve inheritance, polymorphism, or custom copy and move
    semantics.
\end{itemize}

\subsection{References}
\label{c1:sec:references}
References provide an alternative way to pass variables to functions.

A reference acts as an alias for another variable. This means that operations performed on the
reference directly affect the original variable. References behave similarly to pointers but can
be used with normal variable syntax.

References are declared using the \code{&} symbol.

Consider the following example:

\begin{cppcode}
void toggle(bool& state)
{
    state = !state;
}
\end{cppcode}

\noindent The parameter \code{state} is a reference to a \code{bool}.

This means the function modifies the original variable that was passed to it.

Example usage:

\begin{cppcode}
bool state{false};

toggle(state);
\end{cppcode}

\noindent After the function call, the value of \code{state} will be \code{true}.

In C, similar behavior would typically be implemented using pointers:

\begin{ccode}
void toggle(bool* state)
{
    *state = !(*state);
}
\end{ccode}

\subsubsection{Comparison}
\paragraph{C-style approach (with pointer)}

\begin{ccode}
toggle(&state);
\end{ccode}

\paragraph{C++ approach (with reference)}

\begin{cppcode}
toggle(state);
\end{cppcode}

\noindent While pointers are still used in C++, references are often preferred because they:
\begin{itemize}
  \item Cannot be \code{nullptr} (\code{NULL} in C).
  \item Do not require the address operator (\code{&}) to pass addresses.
  \item Do not require dereferencing (\code{*}) to read the value a pointer is pointing at.
  \item Provide clearer syntax.
\end{itemize}
As a result, references are commonly used in modern C++ when a function needs to modify an
existing variable. However, a reference cannot be reseated to refer to another object after
initialization; pointers are still used when we need to change the address we refer to, for
instance when implementing dynamic arrays.

\subsection{The \code{auto} keyword}
\label{c1:sec:auto}
The keyword \code{auto} allows the compiler to automatically deduce the type of a variable from
its initializer.

For example:

\begin{cppcode}
auto number  = 5;   // Deduced as int.
auto counter = 10U; // Deduced as unsigned int.
auto voltage = 3.3; // Deduced as double.
\end{cppcode}

\noindent This can improve readability and reduce the need to repeat long type names.

However, care must be taken when using \code{auto} with brace initialization \code{{}} in
C++11/14. In some cases the compiler may deduce an \code{std::initializer_list} instead of the
intended type.

For example:

\begin{cppcode}
auto a = 10; // Deduced as int.
auto b{10};  // Deduced std::initializer_list<int> prior to C++17.
\end{cppcode}

\noindent Because of this, when using \code{auto} before C++17 it was generally recommended to use
the assignment operator \code{=} instead of brace initialization \code{{}}.

In modern C++ (C++17 and later), the behavior was improved and brace initialization works more
intuitively. However, a common rule of thumb is still:
\begin{itemize}
  \item When using \code{auto}, prefer \code{=} instead of \code{{}} unless you explicitly want an
    \code{std::initializer_list}.
\end{itemize}
However, in this book the data type will usually be specified explicitly. The \code{auto} keyword
will mainly be used in situations where the type would otherwise be verbose, for example when
iterating over containers:

\begin{cppcode}
// Enable all LEDs in a list.
for (auto& led : leds)
{
    led.on();
}
\end{cppcode}

\noindent Without \code{auto}, the type may become very long.

\subsection{Function templates: a first look}
\label{c1:sec:templates}
\note Here we only introduce the basic idea. Templates are covered in more detail in
Chapter~\ref{c5:ch:templates}.

Templates allow functions and classes to operate on multiple data types without duplicating code.

In embedded systems, templates are often used to implement generic utilities such as bit
manipulation helpers.

Consider the following example to set a bit \code{bit} in a given register \code{reg}:
\begin{itemize}
  \item The data type of the register for which to write is set to \code{T}. The type is set at
    compile time when the function is instantiated.
  \item We constrain this function to only work for integral types with a \code{static_assert}. If
    this function is called with some other type, such as a floating point number, a compiler
    error with error message \code{Failed to set bit in register: T must be of integral type!} is
    generated.
  \item We check the data type using \code{std::is_integral<T>} from \code{<type_traits>},
    specifically via its member variable \code{value}:
  \begin{itemize}
    \item If \code{T} is of integral type, \code{std::is_integral<T>::value} is \code{true} and the
      assertion succeeds.
    \item If \code{T} isn't of integral type, \code{std::is_integral<T>::value} is \code{false} and
      the assertion fails.
    \item For more information about type traits, see \secref{c1:sec:typetraits}.
  \end{itemize}
  \item The value \code{1} is cast to type \code{T} via a \code{static_cast}, which is a safer
    alternative to regular C casts.
\end{itemize}

\begin{cppcode}
#include <cstdint>
#include <type_traits>

/**
 * @brief Set bit in register.
 *
 * @tparam T The register type. Must be integral.
 *
 * @param[in, out] reg Register to write to.
 * @param[in] bit The bit to set.
 */
template<typename T>
constexpr void set(T& reg, const std::uint8_t bit) noexcept
{
    static_assert(std::is_integral<T>::value,
        "Failed to set bit in register: T must be of integral type!");
    reg |= (static_cast<T>(1U) << bit);
}
\end{cppcode}

\noindent The function can be used to set bits in registers of arbitrary integral type. In the
example below, bits 1--2 in an 8-bit register are set by calling the \code{set()} function
twice:

\begin{cppcode}
std::uint8_t reg{};
set(reg, 1U);
set(reg, 2U);
\end{cppcode}

\note Unlike C, functions in C++ can have the same name as long as their parameter lists differ.
The compiler determines which version to use at each function call based on the types of the
arguments. This is called function overloading.

\subsubsection{Template instantiation and code size}
Templates are instantiated separately for each data type used in the program. For example, if the
function \code{set()} is used with both \code{std::uint8_t} and \code{std::uint32_t}:

\begin{cppcode}
std::uint8_t reg1{};
std::uint32_t reg2{};

set(reg1, 1U);
set(reg2, 12U);
\end{cppcode}

\noindent Then the compiler generates two separate versions of the function:

\begin{cppcode}
set(std::uint8_t&, std::uint8_t)
set(std::uint32_t&, std::uint8_t)
\end{cppcode}

\noindent Each version is compiled independently and included in the final binary:
\begin{itemize}
  \item This behavior is one of the reasons templates are considered zero-cost abstractions; the
    generated code is specialized for the exact type being used, which allows the compiler to
    optimize the implementation.
  \item However, in embedded systems this also means that excessive template usage can increase
    the binary size, since multiple versions of the same function may be generated.
  \item For this reason, templates should be used carefully in resource-constrained systems,
    especially when many different data types are involved.
\end{itemize}

\subsubsection{Parameter packs}
We can also implement parameter packs so that multiple bits can be set at once:

\begin{cppcode}
#include <cstdint>
#include <type_traits>

template<typename T, typename... Bits>
constexpr void set(T& reg, const Bits... bits) noexcept
{
    static_assert(std::is_integral<T>::value,
        "Failed to set bit in register: T must be of integral type!");

    for (const auto& bit : {bits...})
    {
        reg |= (static_cast<T>(1U) << bit);
    }
}
\end{cppcode}

\noindent In the example below, bits 1--5 in an 8-bit register are set by calling the
\code{set()} function once:

\begin{cppcode}
// Set bit 1-5 in a register.
std::uint8_t reg{};
set(reg, 1U, 2U, 3U, 4U, 5U);
\end{cppcode}

\noindent The compiler will generate code similar to the following:

\begin{cppcode}
reg |= (static_cast<T>(1U) << 1U);
reg |= (static_cast<T>(1U) << 2U);
reg |= (static_cast<T>(1U) << 3U);
reg |= (static_cast<T>(1U) << 4U);
reg |= (static_cast<T>(1U) << 5U);
\end{cppcode}

\subsubsection{Fold expressions (C++17)}
In the implementation above we iterate over the bits using a loop and a so-called \textbf{braced
initializer list}. This approach is easy to read and resembles how similar logic might be written
in C.

Since C++17, we can omit the loop by using a fold expression:

\begin{cppcode}
((reg |= (static_cast<T>(1U) << bits)), ...);
\end{cppcode}

\noindent This expression expands the parameter pack \code{bits...} and performs the operation once
for each element.

\note The operation being folded needs its own pair of parentheses, hence the double parentheses at
the start. The language requires each operand of a fold expression to be a so-called
\emph{cast-expression}, and an assignment such as \code{reg |= x} is not one. Writing

\begin{cppcode}
(reg |= (static_cast<T>(1U) << bits), ...); // Won't compile.
\end{cppcode}

\noindent therefore produces a compiler error such as
\code{binary expression in operand of fold-expression}.

Because of this feature, the \code{set()} function can be implemented more compactly:

\begin{cppcode}
template<typename T, typename... Bits>
constexpr void set(T& reg, const Bits... bits) noexcept
{
    static_assert(std::is_integral<T>::value,
        "Failed to set bit in register: T must be of integral type!");
    ((reg |= (static_cast<T>(1U) << bits)), ...);
}
\end{cppcode}

\subsection{\code{[[nodiscard]]}}
\label{c1:sec:nodiscard}
\code{[[nodiscard]]} is a C++17 attribute that can be added to a function declaration. It tells the
compiler to emit a warning if the caller ignores (discards) the function's return value.

Like \code{constexpr} and templates, \code{[[nodiscard]]} is a zero-cost abstraction: it exists
purely as a compile-time aid and has no effect on the generated machine code or runtime
performance.

Consider the \code{read()} method of the \code{driver::Gpio} struct introduced earlier:

\begin{cppcode}
[[nodiscard]] bool read() const noexcept { return myState; }
\end{cppcode}

\noindent Since \code{read()} has no side effects, calling it and ignoring the result does nothing
useful and is almost always a mistake. With \code{[[nodiscard]]} in place, the following now
produces a compiler warning:

\begin{cppcode}
led.read(); // Warning: ignoring return value of function declared 'nodiscard'.
\end{cppcode}

\noindent While the return value must be used to avoid the warning:

\begin{cppcode}
const bool state{led.read()};
\end{cppcode}

\noindent \code{[[nodiscard]]} is particularly useful in embedded systems for functions such as:
\begin{itemize}
  \item Query/read functions, similar to \code{read()} above.
  \item Functions returning error or status codes, where ignoring the result may hide a failure.
  \item Factory functions that return an owning resource, such as \code{std::unique_ptr}
    (introduced in Chapter~\ref{c4:ch:factory}), where discarding the result would immediately
    release the resource.
\end{itemize}
\note \code{[[nodiscard]]} should not be added indiscriminately. Functions that are called mainly
for their side effects, where the return value is genuinely optional, do not need it.

\subsection{Additional information about type traits}
\label{c1:sec:typetraits}
Type traits allow programs to inspect and reason about types at compile time.

In the C++ standard library, many type traits are implemented as small structs containing a static
constant boolean value (\code{static} means that the value belongs to the type itself rather than
to an instance of the struct).

We can implement a simplified version of a trait that checks whether a type is unsigned:

\begin{cppcode}
template<typename T>
struct isUnsigned
{
    static constexpr bool value{false};
};
\end{cppcode}

\note
\begin{itemize}
  \item The type trait is nothing more than a struct template containing a boolean constant
    \code{value}.
  \item By default, this boolean value is \code{false} for all types \code{T}.
  \item We can specialize the template for specific types \code{T} so that \code{value} becomes
    \code{true}, as shown below, where the types \code{std::uint8_t}, \code{std::uint16_t},
    \code{std::uint32_t}, and \code{std::uint64_t} are treated as unsigned types:
\end{itemize}

\begin{cppcode}
template<>
struct isUnsigned<std::uint8_t>
{
    static constexpr bool value{true};
};

template<>
struct isUnsigned<std::uint16_t>
{
    static constexpr bool value{true};
};

template<>
struct isUnsigned<std::uint32_t>
{
    static constexpr bool value{true};
};

template<>
struct isUnsigned<std::uint64_t>
{
    static constexpr bool value{true};
};
\end{cppcode}

\note There is no separate specialization for \code{std::size_t}. It is not a type of its own but
an alias of one of the unsigned types above (\code{std::uint64_t} on a 64-bit Linux machine,
typically \code{std::uint32_t} on a 32-bit microcontroller), so it is already covered, and
specializing it as well would be a redefinition that fails to compile.


\noindent We can use our own type trait \code{isUnsigned<T>} to check if a type \code{T} is
unsigned as shown below:

\begin{cppcode}
static_assert(isUnsigned<T>::value, "T must be of unsigned type!");
\end{cppcode}

\noindent Since C++17, it is common to also provide a variable template ending in \code{_v} that
exposes the boolean value directly:

\begin{cppcode}
template<typename T>
inline constexpr bool isUnsigned_v{isUnsigned<T>::value};
\end{cppcode}

\note \code{inline} allows the variable template to be defined in a header file and included in
multiple translation units without causing multiple definition errors.

This variable template allows us to write:

\begin{cppcode}
static_assert(isUnsigned_v<T>, "T must be of unsigned type!");
\end{cppcode}

\noindent In this book, the traditional form \code{isUnsigned<T>::value} will be used for clarity.

\subsection{Summary}
\label{c1:sec:cppsummary}
This section introduced the following modern C++ features commonly used in embedded software:
\begin{itemize}
  \item \code{.h} vs \code{.hpp} header file conventions.
  \item \code{extern "C"} and \code{#ifdef __cplusplus} for headers shared between C and C++.
  \item Namespaces.
  \item \code{constexpr} and compile-time constants.
  \item \code{noexcept} and exception handling in embedded systems.
  \item Default arguments.
  \item Structs with member functions, constructors/destructors, and encapsulation.
  \item References.
  \item The \code{auto} keyword.
  \item Function templates.
  \item \code{[[nodiscard]]} for catching ignored return values.
\end{itemize}
These features allow developers to write safer, clearer, and more maintainable embedded software
while still maintaining full control over hardware and performance.
